\documentclass[11pt]{article}

\usepackage[margin=1.25in]{geometry}
\usepackage[T1]{fontenc}
\usepackage{amsmath,amssymb,amsthm,mathtools}
\usepackage{bm}
\usepackage{enumitem}
\usepackage{microtype}
\usepackage{xcolor}
\usepackage[colorlinks=true,linkcolor=blue!55!black,citecolor=blue!55!black,urlcolor=blue!55!black]{hyperref}

\newtheorem{theorem}{Theorem}[section]
\newtheorem{lemma}[theorem]{Lemma}
\newtheorem{proposition}[theorem]{Proposition}
\newtheorem{corollary}[theorem]{Corollary}
\newtheorem{definition}[theorem]{Definition}
\theoremstyle{remark}
\newtheorem{remark}[theorem]{Remark}

\newcommand{\R}{\mathbb{R}}
\newcommand{\E}{\mathbb{E}}
\newcommand{\Vol}{\operatorname{Vol}}
\newcommand{\supp}{\operatorname{supp}}
\newcommand{\eps}{\varepsilon}
\newcommand{\one}{\bm{1}}
\newcommand{\ip}[2]{\left\langle #1,#2\right\rangle}
\newcommand{\MinkDiff}{\mathbin{\ominus}}

\title{Linear-size sparsifiers for sums of seminorms}
\author{}
\date{July 12, 2026}

\begin{document}

\maketitle

\begin{abstract}
Let $N_1,\ldots,N_m$ be seminorms on $\R^n$, and write
$N=\sum_{i=1}^m N_i$.  We prove that, for every
$\eps\in(0,1/2]$, there are nonnegative weights $w_1,\ldots,w_m$, at
most
\[
  O\!\left(\frac{d}{\eps^2}\log\frac1\eps\right)
\]
of which are nonzero, such that
\[
  (1-\eps)N(x)\leq \sum_{i=1}^m w_iN_i(x)
  \leq (1+\eps)N(x)\qquad\text{for every }x\in\R^n.
\]
Here $d=\dim(\R^n/\ker N)\leq n$ is the effective dimension of the
sum.  This removes the dimension-dependent polylogarithmic losses in
the previous general bounds for sums of seminorms.  As consequences,
sums of normalized symmetric submodular functions admit sparsifiers
of the same size, with the natural effective dimension, and every
weighted hypergraph on $n$ vertices has a cut sparsifier with
$O((n-\kappa)\eps^{-2}\log(1/\eps))$ hyperedges, where $\kappa$ is the
number of connected components of its positive-weight $2$-section.

At a high level, we represent seminorms as support functions of
compact centrally symmetric convex sets and extend the recent
coefficient-space argument of Reis and Rothvoss from sums of segments
to arbitrary Minkowski sums.  Convexity of volume under Minkowski
erosion supplies the required volumetric estimate, and partial
coloring then eliminates a constant fraction of the summands at each
step.
\end{abstract}

\begin{quote}
\small\noindent\textbf{Note.} This paper was fully AI generated by GPT
5.6 Sol and the Devin harness, prompted by Yang Liu.
\end{quote}

\section{Introduction}

Consider seminorms $N_1,\ldots,N_m:\R^n\to\R_+$ and their sum
\begin{equation}\label{eq:sum-norm}
  N(x)=\sum_{i=1}^m N_i(x).
\end{equation}
The sparsification problem asks whether most of the summands can be
discarded after reweighting while preserving $N$ simultaneously on all
of $\R^n$.  More precisely, one seeks a nonnegative vector
$w\in\R_+^m$ with small support for which
\begin{equation}\label{eq:approximation}
  \left|N(x)-\sum_{i=1}^m w_iN_i(x)\right|
  \leq \eps N(x)\qquad\text{for every }x\in\R^n.
\end{equation}
The familiar special case $N_i(x)=|\ip{a_i}{x}|$ is row
sparsification in $\ell_1$.  Graph cut sparsification is obtained by
taking $a_i$ to be signed edge-incidence vectors.

Jambulapati, Lee, Liu, and Sidford~\cite{JLLS23} proved that arbitrary
sums of seminorms admit sparsifiers with
\[
  O\!\left(\eps^{-2}n\log(n/\eps)(\log n)^{5/2}\right)
\]
nonzero weights.  Their result also gives an efficient construction
under a standard rounding assumption.  For one-dimensional summands,
Talagrand's classical theorem~\cite{Talagrand90} gives
$O(\eps^{-2}n\log n)$ rows.  Very recently, Reis and
Rothvoss~\cite{RR26} removed the dimension-dependent logarithm and
proved the bound $O(\eps^{-2}n\log(1/\eps))$ for this case.

Our main result gives the latter bound for arbitrary seminorms.  It is
convenient to state it in terms of the dimension on which the sum is a
genuine norm.

\begin{theorem}\label{thm:main}
Let $N_1,\ldots,N_m$ be seminorms on $\R^n$, let
$N=\sum_{i=1}^mN_i$, and set
\[
  d=\dim(\R^n/\ker N).
\]
For every $\eps\in(0,1/2]$, there is a vector $w\in\R_+^m$ such that
\[
  |\supp(w)|\leq
  C\frac{d}{\eps^2}\log\frac1\eps
\]
and
\[
  (1-\eps)N(x)\leq\sum_{i=1}^m w_iN_i(x)
  \leq(1+\eps)N(x)
  \qquad\text{for every }x\in\R^n.
\]
Here $C>0$ is a universal constant.
\end{theorem}

The dependence on $d$ is necessary already for graph cuts, up to the
$\log(1/\eps)$ factor; see, for example, the lower bounds discussed
in~\cite{AKW14}.  Theorem~\ref{thm:main} is existential.  The
Reis--Rothvoss framework is also existential in the $\ell_1$ setting
and does not yield a polynomial-time construction~\cite{RR26}.

\paragraph{Proof idea.}
Every seminorm is the support function of a compact centrally
symmetric convex set.  Thus, if $C_i$ is the support body of $N_i$,
then $N$ is the support function of the Minkowski sum
$C=C_1+\cdots+C_m$.  Condition~\eqref{eq:approximation} is equivalent
to
\[
  (1-\eps)C\subseteq \sum_{i=1}^m w_iC_i
  \subseteq(1+\eps)C.
\]
Reis and Rothvoss prove this inclusion when every $C_i$ is a segment.
Their proof has two parts.  A volumetric argument shows that the set of
small simultaneous perturbations of the coefficients has large
volume, and a partial-coloring theorem then finds a perturbation that
sets a constant fraction of the coefficients to zero.  Iterating the
step gives the desired sparsity.

Only one part of the volumetric argument uses that the summands are
segments: the assertion that volume is separately convex under
Minkowski erosion by the summands.  We prove the following general
replacement:
\begin{equation}\label{eq:erosion-convex-intro}
  t\longmapsto \Vol(K\MinkDiff tL)
  \quad\text{is convex on }\R_+
\end{equation}
for arbitrary convex bodies $K,L$.  For full-dimensional $L$, the
derivative of this function is minus the anisotropic perimeter of the
eroded body.  That perimeter is a mixed volume, hence decreases under
set inclusion.  Lower-dimensional $L$ follow by approximation.  This
one-variable statement implies separate convexity when the summands
are arbitrary support bodies, after which the coefficient-space proof
of~\cite{RR26} applies without change.  We include all steps needed for
the reduction.

\paragraph{Applications.}
The Lov\'asz extension of a nonnegative normalized symmetric
submodular function is a seminorm~\cite{Lovasz83}.  We therefore obtain
a sparsifier of size
$O(n\eps^{-2}\log(1/\eps))$ for every decomposable symmetric
submodular function, improving the earlier general bound
from~\cite{JLLS23}.  A hyperedge cut is one such
function.  Equivalently, its Lov\'asz extension is the range seminorm
on the hyperedge.  This yields a weighted hypergraph cut sparsifier
with $O(n\eps^{-2}\log(1/\eps))$ hyperedges.  Prior polynomial-time
constructions use $O(n\eps^{-2}\log n)$ hyperedges~\cite{CKN20}; our
conclusion concerns existence rather than efficient construction.

\paragraph{Organization.}
Section~\ref{sec:preliminaries} records the convex-geometric
translation and the coefficient-space results from~\cite{RR26}.
Section~\ref{sec:erosion} proves~\eqref{eq:erosion-convex-intro} and a
random inclusion theorem for Minkowski sums.  In
Section~\ref{sec:proof}, we establish the large-volume perturbation
body and prove Theorem~\ref{thm:main}.  Section~\ref{sec:applications}
gives the submodular and hypergraph consequences.

\section{Preliminaries}\label{sec:preliminaries}

All convex sets in this paper are subsets of a finite-dimensional
Euclidean space.  A \emph{convex body} is a compact convex set with
nonempty interior in the ambient space.  We allow other compact convex
sets to be lower-dimensional.  For a compact convex set $K$, its
support function is
\[
  h_K(x)=\max_{y\in K}\ip{y}{x}.
\]
Support functions characterize compact convex sets, and satisfy
$h_{K+L}=h_K+h_L$ and $h_{tK}=t h_K$ for $t\geq0$.  In particular,
$K\subseteq L$ if and only if $h_K\leq h_L$ pointwise.

For compact convex sets $K,L$, with $L$ nonempty, define their
Minkowski erosion (also called Minkowski subtraction) by
\[
  K\MinkDiff L=\{x:x+L\subseteq K\}.
\]
This is a compact convex set, possibly empty, and
\begin{equation}\label{eq:erosion-associative}
  (K\MinkDiff L)\MinkDiff M=K\MinkDiff(L+M).
\end{equation}
If $K$ is convex and $0\leq t\leq1$, then
\begin{equation}\label{eq:self-erosion}
  K\MinkDiff tK=(1-t)K.
\end{equation}
We use $\Vol_k$ for $k$-dimensional Lebesgue measure, with the
convention $\Vol_0(\{0\})=1$.

\subsection{Support bodies of seminorms}

Let $N$ be a seminorm on $\R^n$.  Its kernel is a linear subspace.  If
$E=(\ker N)^\perp$, then the restriction of $N$ to $E$ is a norm and
$N(x)=N(P_Ex)$, where $P_E$ is orthogonal projection onto $E$.

\begin{lemma}\label{lem:support-body}
For every seminorm $N$ on $\R^n$, the set
\[
  C_N=\{y\in\R^n:\ip{y}{x}\leq N(x)\text{ for every }x\in\R^n\}
\]
is compact, convex, centrally symmetric, and contained in
$(\ker N)^\perp$.  Moreover,
\[
  N(x)=h_{C_N}(x)\qquad\text{for every }x\in\R^n.
\]
If $N$ is a norm on a Euclidean space $E$, then $C_N$ has nonempty
interior in $E$.
\end{lemma}

\begin{proof}
Convexity, closedness, and central symmetry follow directly from the
definition.  If $z\in\ker N$ and $y\in C_N$, then
$|\ip{y}{z}|\leq N(z)=0$, so $C_N\subseteq(\ker N)^\perp$.  On that
orthogonal complement, equivalence of finite-dimensional norms gives
$N(x)\leq R\|x\|_2$ for some $R<\infty$.  Testing the defining
inequality in the direction $x=y$ shows $\|y\|_2\leq R$ for
$y\in C_N$.  Hence $C_N$ is compact.

The inequality $h_{C_N}\leq N$ is part of the definition.  For the
reverse inequality, fix $x$.  If $x\in\ker N$, both sides vanish.
Otherwise, apply Hahn--Banach to the linear functional on
$\operatorname{span}\{x\}$ that maps $tx$ to $tN(x)$, dominated by
$N$.  Its extension has the form $z\mapsto\ip{y}{z}$ for some
$y\in C_N$ and satisfies $\ip{y}{x}=N(x)$.  Thus $h_{C_N}(x)\geq
N(x)$.  Finally, when $N$ is a norm, norm equivalence also gives
$N(x)\geq r\|x\|_2$ for some $r>0$, and therefore
$rB_2\subseteq C_N$.
\end{proof}

For seminorms $N_1,\ldots,N_m$, let $C_i=C_{N_i}$.  With
$N=\sum_iN_i$, we have
\begin{equation}\label{eq:support-sum}
  N=h_C,\qquad C=\sum_{i=1}^m C_i.
\end{equation}
If $E=(\ker N)^\perp$, then every $C_i\subseteq E$, and $C$ is a
centrally symmetric convex body in the $d$-dimensional space $E$.
Indeed, $\ker N=\bigcap_i\ker N_i$ because the summands are
nonnegative.

\subsection{Two coefficient-space theorems}

We record the two results from Reis and Rothvoss that drive the support
reduction.  For $A\subseteq\R^m$ and $S\subseteq[m]$, write
\[
  A_S=\{x_S\in\R^S:(x_S,0_{[m]\setminus S})\in A\}
\]
for its coordinate section.

\begin{theorem}[Reis--Rothvoss symmetrizer theorem]\label{thm:RR-symmetrizer}
Let $p,\eta\in(0,1/2]$, and suppose
$m\geq\log_2(1/p)/\eta^2$.  Let $A\subseteq[-1,1]^m$ be a convex body
such that $[-\eta,\eta]^m\subseteq A$ and
\[
  \Vol_{|S|}(A_S)\geq p\,2^{|S|}
  \qquad\text{for every }S\subseteq[m].
\]
Then
\[
  \Vol_m(A\cap(-A))\geq2^{-5m}.
\]
\end{theorem}

This is Theorem~16 of~\cite{RR26}.  We also use their partial-coloring
statement, which is a volume formulation of a result developed
in~\cite{RR23}.

\begin{theorem}[Reis--Rothvoss partial coloring]\label{thm:RR-coloring}
For every $c>0$, there is a constant $s=s(c)>0$ with the following
property.  If $Q\subseteq[-1,1]^m$ is a centrally symmetric convex
body and $\Vol_m(Q)\geq c^m$, then there is
$z\in sQ\cap[-1,1]^m$ such that
\[
  |\{i\in[m]:|z_i|=1\}|\geq m/2.
\]
\end{theorem}

This is Theorem~7 of~\cite{RR26}.  Only the existence of a universal
constant $s=s(2^{-5})$ will be relevant below.

\section{Minkowski erosion and random sums}\label{sec:erosion}

The main new geometric input is the following extension of the segment
erosion lemma used in~\cite{RR26}.

\begin{lemma}[Convexity of erosion volume]\label{lem:erosion-convex}
Let $K,L\subseteq\R^d$ be compact convex sets, with $L$ nonempty.  Set
$\Vol_d(\varnothing)=0$.  Then
\[
  f(t)=\Vol_d(K\MinkDiff tL),\qquad t\geq0,
\]
is convex.
\end{lemma}

\begin{proof}
If $K$ is lower-dimensional, then $f\equiv0$, so assume that $K$ is
full-dimensional.  We first suppose that $L$ is full-dimensional.
Translating $L$ by a vector $a$ translates
$K\MinkDiff tL$ by $-ta$ and does not change its volume, so we may
assume $0\in\operatorname{int}(L)$.  Let
\[
  r=r(K;L)=\sup\{t\geq0:K\MinkDiff tL\neq\varnothing\}
\]
be the relative inradius, and write $K_t=K\MinkDiff tL$ for
$0\leq t\leq r$.  For $t<r$, the set $K_t$ is full-dimensional.

Recall that the anisotropic perimeter of a convex body $A$ relative to
$L$ is
\[
  P_L(A)=d\,V(A[d-1],L),
\]
where $V$ denotes mixed volume.  The inner-parallel-body first
variation formula gives
\begin{equation}\label{eq:first-variation}
  \frac{d}{dt}\Vol_d(K_t)=-P_L(K_t),\qquad 0<t<r.
\end{equation}
See, for example, Proposition~2.6 of~\cite{RS22} or
Section~7.5 of~\cite{Schneider14}.  If $0\leq t<u<r$, then
$K_u\subseteq K_t$.  Monotonicity of mixed volumes in each argument
therefore gives
\[
  P_L(K_u)\leq P_L(K_t).
\]
Consequently, the derivative in~\eqref{eq:first-variation} is
nondecreasing, and $f$ is convex on $[0,r)$.  At $t=r$, $K_r$ has
empty interior and hence zero $d$-dimensional volume; otherwise a
slightly larger translate of $L$ would fit in $K$.  Extending $f$ by
zero on $[r,\infty)$ preserves convexity because its left derivatives
are nonpositive and the right derivative after $r$ is zero.

It remains to remove the full-dimensionality assumption on $L$.  Let
$B$ be the Euclidean unit ball and $L_\delta=L+\delta B$ for
$\delta>0$.  The functions
\[
  f_\delta(t)=\Vol_d(K\MinkDiff tL_\delta)
\]
are convex by the first part.  By associativity of erosion,
\begin{equation}\label{eq:regularized-erosion}
  K\MinkDiff tL_\delta
  =(K\MinkDiff tL)\MinkDiff t\delta B.
\end{equation}
As $\delta\downarrow0$, the sets on the right increase to the interior
of $K\MinkDiff tL$ when that set is full-dimensional; their volumes
converge to its volume.  If $K\MinkDiff tL$ is lower-dimensional, both
sides of the desired volume limit are zero.  Hence
$f_\delta(t)\to f(t)$ pointwise for every $t\geq0$.  A pointwise limit
of finite convex functions is convex, which completes the proof.
\end{proof}

\begin{corollary}[Separate convexity]\label{cor:separate-convex}
Let $K,L_1,\ldots,L_m\subseteq\R^d$ be compact convex sets, with each
$L_i$ nonempty.  The function
\[
  F(t)=\Vol_d\!\left(K\MinkDiff\sum_{i=1}^m t_iL_i\right),
  \qquad t\in\R_+^m,
\]
is convex in each coordinate separately.
\end{corollary}

\begin{proof}
Fix all coordinates except $t_i$ and use
associativity~\eqref{eq:erosion-associative} to write the displayed
erosion as $K'\MinkDiff t_iL_i$ for a fixed compact convex set $K'$.
If $K'$ is empty, the function is zero.  Otherwise,
Lemma~\ref{lem:erosion-convex} applies.
\end{proof}

We next generalize the random-zonotope inclusion theorem
from~\cite{RR26} to arbitrary Minkowski summands.

\begin{theorem}[Random Minkowski-sum inclusion]\label{thm:random-inclusion}
Let $K\subseteq\R^d$ be a centrally symmetric convex body, and let
$L_1,\ldots,L_m\subseteq\R^d$ be nonempty compact centrally symmetric
convex sets.  For $t\in\R_+^m$, write
$L_t=\sum_{i=1}^m t_iL_i$.  If $X\in\R_+^m$ has independent
integrable coordinates, then
\[
  \Pr[L_X\subseteq K]
  \geq
  \frac{\Vol_d(K\MinkDiff L_{\E X})}{\Vol_d(K)}.
\]
\end{theorem}

\begin{proof}
Set $F(t)=\Vol_d(K\MinkDiff L_t)$.  By
Corollary~\ref{cor:separate-convex}, $F$ is separately convex.  Repeated
one-dimensional Jensen inequalities, using independence of the
coordinates, give
\begin{equation}\label{eq:separate-jensen}
  \E F(X)\geq F(\E X).
\end{equation}
Every erosion $K\MinkDiff L_X$ is contained in $K$ and therefore has
volume at most $\Vol_d(K)$.  Moreover, positive volume implies that the
erosion is nonempty, which for centrally symmetric $K$ and $L_X$ is
equivalent to $L_X\subseteq K$.  It follows that
\[
  \Pr[L_X\subseteq K]\Vol_d(K)
  \geq \E F(X)
  \geq F(\E X),
\]
where the last inequality is~\eqref{eq:separate-jensen}.
\end{proof}

\section{Sparsification}\label{sec:proof}

We first show that the body of coefficient perturbations preserving a
Minkowski sum has large volume.  This is the point at which
Theorem~\ref{thm:random-inclusion} enters the Reis--Rothvoss proof.

Let $d\geq1$.  Let $C_1,\ldots,C_m$ be nonempty compact centrally
symmetric convex sets in $\R^d$, and suppose
$C=\sum_{i=1}^m C_i$ is full-dimensional.  For $t\in\R_+^m$, write
$C_t=\sum_i t_iC_i$.

\begin{proposition}[Large perturbation body]\label{prop:large-body}
Fix $\eta\in(0,1/2]$, and define
\[
  A_\eta=
  \{z\in[-1,1]^m:C_{\one+z}\subseteq(1+\eta)C\}.
\]
Then $A_\eta$ is convex, contains $[-\eta,\eta]^m$, and for every
$S\subseteq[m]$,
\[
  \Vol_{|S|}((A_\eta)_S)
  \geq \left(\frac{\eta}{2}\right)^d2^{|S|}.
\]
Consequently, if
\begin{equation}\label{eq:large-body-threshold}
  m\geq \frac{2d\log_2(1/\eta)}{\eta^2},
\end{equation}
then the centrally symmetric convex body
\[
  Q_\eta=A_\eta\cap(-A_\eta)
  =\{z\in[-1,1]^m:(1-\eta)C\subseteq C_{\one+z}
      \subseteq(1+\eta)C\}
\]
satisfies $\Vol_m(Q_\eta)\geq2^{-5m}$.
\end{proposition}

\begin{proof}
Taking support functions gives
\begin{equation}\label{eq:coefficient-halfspaces}
  A_\eta=
  \left\{z\in[-1,1]^m:
  \sum_{i=1}^m z_i h_{C_i}(x)
  \leq \eta\sum_{i=1}^m h_{C_i}(x)
  \text{ for every }x\in\R^d\right\}.
\end{equation}
Thus $A_\eta$ is an intersection of halfspaces and is convex.  The
same display shows $[-\eta,\eta]^m\subseteq A_\eta$.

Fix $S\subseteq[m]$.  Let $X_i$ be uniform on $[0,2]$ for $i\in S$
and set $X_i=1$ deterministically for $i\notin S$.  The coordinates are
independent and $\E X=\one$.  Theorem~\ref{thm:random-inclusion},
applied with target $(1+\eta)C$, yields
\begin{align*}
  \frac{\Vol_{|S|}((A_\eta)_S)}{2^{|S|}}
  &=\Pr[C_X\subseteq(1+\eta)C]\\
  &\geq
  \frac{\Vol_d((1+\eta)C\MinkDiff C)}
       {\Vol_d((1+\eta)C)}
   =\left(\frac{\eta}{1+\eta}\right)^d
   \geq\left(\frac{\eta}{2}\right)^d,
\end{align*}
where~\eqref{eq:self-erosion} gives the equality.

Intersecting~\eqref{eq:coefficient-halfspaces} with its negative shows
that $Q_\eta=A_\eta\cap(-A_\eta)$ has the claimed two-sided inclusion
description.  Set $p=(\eta/2)^d$.  Since $\eta\leq1/2$,
\[
  \log_2(1/p)=d\log_2(2/\eta)
  \leq2d\log_2(1/\eta).
\]
Under~\eqref{eq:large-body-threshold}, all hypotheses of
Theorem~\ref{thm:RR-symmetrizer} hold, and hence
$\Vol_m(Q_\eta)\geq2^{-5m}$.
\end{proof}

We now iterate a partial coloring.  The following elementary estimate
will control the accumulated error.

\begin{lemma}\label{lem:phi}
Let $\phi(u)=\log_2(1/u)/u^2$ for $u\in(0,1/2]$.  For every
$a\geq1$ and $u\in(0,1/2]$,
\[
  a\phi(u)\geq\phi(ua^{-1/4}).
\]
\end{lemma}

\begin{proof}
Writing $q=\log_2a\geq0$ gives
\[
  \frac{\phi(ua^{-1/4})}{\phi(u)}
  =a^{1/2}\left(1+\frac{q}{4\log_2(1/u)}\right)
  \leq a^{1/2}\left(1+\frac q4\right)
  \leq a.
\]
The last inequality is $1+q/4\leq2^{q/2}$, which follows from
convexity of the right-hand side and its derivative at $q=0$.
\end{proof}

\begin{proof}[Proof of Theorem~\ref{thm:main}]
If $d=0$, then $N\equiv0$ and the zero vector of weights suffices.
Otherwise, work in the $d$-dimensional space
$E=(\ker N)^\perp$.  Let $C_i$ be the support body of $N_i$ in $E$,
as in Lemma~\ref{lem:support-body}, so that
$C=\sum_iC_i$ is full-dimensional and~\eqref{eq:support-sum} holds.
Zero summands may be discarded at the outset.

Let $s=s(2^{-5})$ be the constant in
Theorem~\ref{thm:RR-coloring}.  Choose a sufficiently large universal
constant $C_0$, depending only on $s$, and set
\[
  M=C_0d\phi(\eps)
  =C_0\frac{d}{\eps^2}\log_2\frac1\eps.
\]
If $m\leq M$, take every weight equal to one.  Suppose $m>M$.

Starting from $w^{(0)}=\one$, we construct nonnegative weight vectors
$w^{(t)}$.  Write
\[
  S_t=\supp(w^{(t)}),\qquad m_t=|S_t|,
  \qquad C^{(t)}=\sum_{i\in S_t}w_i^{(t)}C_i.
\]
As long as $m_t>M$, let $\eta_t\in(0,1/2]$ be the unique solution of
\begin{equation}\label{eq:eta-definition}
  m_t=2d\phi(\eta_t).
\end{equation}
Indeed, $\phi$ is continuous and strictly decreasing from infinity to
$4$ on $(0,1/2]$.  Moreover, $m_t>M$ and Lemma~\ref{lem:phi}, applied
with $a=C_0/2$, give
\[
  \eta_t\leq\eps(C_0/2)^{-1/4}.
\]
Thus the choice of $C_0$ ensures existence, $\eta_t\leq\eps$, and
$s\eta_t<1/2$.  Apply
Proposition~\ref{prop:large-body} to the active summands
$w_i^{(t)}C_i$, $i\in S_t$.  It gives a symmetric convex body
\[
  Q_t=\left\{z\in[-1,1]^{S_t}:
  (1-\eta_t)C^{(t)}
  \subseteq\sum_{i\in S_t}(1+z_i)w_i^{(t)}C_i
  \subseteq(1+\eta_t)C^{(t)}\right\}
\]
with $\Vol_{m_t}(Q_t)\geq2^{-5m_t}$.
Theorem~\ref{thm:RR-coloring} supplies
$z^{(t)}\in sQ_t\cap[-1,1]^{S_t}$ for which at least $m_t/2$
coordinates have absolute value one.  Since $Q_t$ is symmetric, we
may replace $z^{(t)}$ by its negative and arrange that at least
$m_t/4$ coordinates equal $-1$.  Update
\[
  w_i^{(t+1)}=(1+z_i^{(t)})w_i^{(t)}
  \quad(i\in S_t),
\]
and keep all other weights zero.  The weights remain nonnegative and
\begin{equation}\label{eq:support-shrink}
  m_{t+1}\leq\frac34m_t.
\end{equation}
Moreover, because $z^{(t)}\in sQ_t$, the support-function description
of $Q_t$ gives
\begin{equation}\label{eq:one-step-error}
  (1-s\eta_t)C^{(t)}\subseteq C^{(t+1)}
  \subseteq(1+s\eta_t)C^{(t)}.
\end{equation}

Let $T$ be the first index for which $m_T\leq M$.  It remains to bound
the accumulated error.  From~\eqref{eq:support-shrink} and
$m_{T-1}>M$, for $0\leq t<T$ we have
\[
  m_t>M\left(\frac43\right)^{T-1-t}.
\]
Combining this with~\eqref{eq:eta-definition} and the definition of
$M$ gives
\[
  \phi(\eta_t)>
  \frac{C_0}{2}\phi(\eps)
  \left(\frac43\right)^{T-1-t}.
\]
Lemma~\ref{lem:phi} and monotonicity of $\phi$ imply
\[
  \eta_t\leq
  \eps\left(\frac{C_0}{2}\right)^{-1/4}
  \left(\frac34\right)^{(T-1-t)/4}.
\]
We choose $C_0$ large enough that
\begin{equation}\label{eq:error-sum}
  \sum_{t=0}^{T-1}s\eta_t
  \leq
  s\eps\left(\frac{C_0}{2}\right)^{-1/4}
  \sum_{j=0}^{\infty}\left(\frac34\right)^{j/4}
  \leq\frac\eps2.
\end{equation}
Iterating~\eqref{eq:one-step-error} and using
\eqref{eq:error-sum},
\begin{align*}
  \prod_{t<T}(1-s\eta_t)
  &\geq1-\sum_{t<T}s\eta_t\geq1-\eps,\\
  \prod_{t<T}(1+s\eta_t)
  &\leq\exp\left(\sum_{t<T}s\eta_t\right)
    \leq e^{\eps/2}\leq1+\eps.
\end{align*}
The final inequality holds for $0\leq\eps\leq1/2$.  We conclude that
\[
  (1-\eps)C\subseteq C^{(T)}\subseteq(1+\eps)C.
\]
Taking support functions and using~\eqref{eq:support-sum} proves the
uniform seminorm approximation.  Finally,
$|\supp(w^{(T)})|=m_T\leq M$, as required.
\end{proof}

\begin{remark}
The proof uses no description or evaluation oracle for the summands.
It therefore establishes existence for completely arbitrary
finite-dimensional seminorms.  Finding the partial coloring requires
a separation oracle for a body defined by uniform comparison of two
sums of seminorms.  We do not assert that such separation, or the
resulting sparsifier, can be computed in polynomial time in a general
oracle model.
\end{remark}

\section{Symmetric submodular functions and hypergraph cuts}
\label{sec:applications}

Let $V$ be a finite ground set of size $n$.  A function
$f:2^V\to\R_+$ is \emph{submodular} if
\[
  f(A)+f(B)\geq f(A\cup B)+f(A\cap B)
  \qquad(A,B\subseteq V),
\]
and \emph{symmetric} if $f(S)=f(V\setminus S)$ for every $S\subseteq
V$.  We call $f$ normalized when $f(\varnothing)=0$; symmetry then
also gives $f(V)=0$.

For a normalized symmetric function, its Lov\'asz extension may be
written, up to immaterial choices at threshold ties, as
\begin{equation}\label{eq:lovasz-extension}
  \bar f(x)=\int_{-\infty}^{\infty}
  f(\{v\in V:x_v\geq t\})\,dt.
\end{equation}
The integral is finite because the threshold set is $V$ below the
minimum coordinate and empty above the maximum coordinate.

\begin{lemma}\label{lem:lovasz-seminorm}
If $f:2^V\to\R_+$ is normalized, symmetric, and submodular, then
$\bar f$ is a seminorm on $\R^V$.  It vanishes on constant vectors and
satisfies
\[
  \bar f(\one_S)=f(S)\qquad(S\subseteq V).
\]
Moreover, the Lov\'asz extension is linear in the set function.
\end{lemma}

\begin{proof}
The Lov\'asz extension of a submodular function is convex
\cite{Lovasz83}.  Formula~\eqref{eq:lovasz-extension} shows positive
homogeneity.  Symmetry of $f$, followed by the change of variables
$t\mapsto-t$, gives $\bar f(-x)=\bar f(x)$; threshold ties affect only
a measure-zero set of $t$.  A convex, positively homogeneous, even
function is a seminorm.  The remaining assertions follow directly
from~\eqref{eq:lovasz-extension}.
\end{proof}

\begin{corollary}[Symmetric submodular sparsification]
\label{cor:submodular}
Let $f_1,\ldots,f_m:2^V\to\R_+$ be normalized symmetric submodular
functions, set $F=\sum_i f_i$, and let
\[
  d_F=n-\dim\left(\bigcap_{i=1}^m\ker\bar f_i\right)\leq n-1.
\]
For every $\eps\in(0,1/2]$, there are nonnegative weights
$w_1,\ldots,w_m$ with
\[
  |\supp(w)|=O\!\left(\frac{d_F}{\eps^2}
  \log\frac1\eps\right)
\]
such that
\[
  (1-\eps)F(S)\leq\sum_{i=1}^m w_if_i(S)
  \leq(1+\eps)F(S)
  \qquad\text{for every }S\subseteq V.
\]
In fact, the same weights approximate the Lov\'asz extension of $F$
on every vector in $\R^V$.
\end{corollary}

\begin{proof}
Apply Theorem~\ref{thm:main} to the seminorms $\bar f_i$ from
Lemma~\ref{lem:lovasz-seminorm}.  Every $\bar f_i$ vanishes on the
constant vectors, so $d_F\leq n-1$.  Evaluate the resulting inequality
at $\one_S$.
\end{proof}

The effective dimension in Corollary~\ref{cor:submodular} can be
smaller than $n-1$ if the common kernel of the Lov\'asz extensions is
larger than the constants.  The conclusion applies, for example, to
sums of graph connectivity functions, hypergraph cut functions, and
matroid connectivity functions.

We spell out the hypergraph case.  Let $H=(V,\mathcal E,c)$ be a
weighted hypergraph, where $c_e\geq0$ for $e\in\mathcal E$.  Its cut
function is
\[
  \delta_H(S)=
  \sum_{\substack{e\in\mathcal E:\
  e\cap S\neq\varnothing,\ e\setminus S\neq\varnothing}}c_e.
\]
For each hyperedge $e$, define the range function
\[
  R_e(x)=\max_{u\in e}x_u-\min_{v\in e}x_v
\]
and the weighted range seminorm
\begin{equation}\label{eq:range-seminorm}
  N_e(x)=c_eR_e(x).
\end{equation}
It is the Lov\'asz extension of the cut indicator of $e$, and
$N_e(\one_S)=c_e$ exactly when $e$ crosses $S$.

\begin{corollary}[Hypergraph cut sparsification]
\label{cor:hypergraph}
For every weighted hypergraph $H=(V,\mathcal E,c)$ on $n$ vertices,
let $\kappa(H)$ be the number of connected components, including
isolated vertices, of the graph that joins $u,v\in V$ when some
positive-weight hyperedge contains both.  For every
$\eps\in(0,1/2]$, there is a weighted subhypergraph
$H'=(V,\mathcal E',c')$ with
\[
  |\mathcal E'|=O\!\left(\frac{n-\kappa(H)}{\eps^2}
  \log\frac1\eps\right)
\]
such that
\[
  (1-\eps)\delta_H(S)\leq\delta_{H'}(S)
  \leq(1+\eps)\delta_H(S)
  \qquad\text{for every }S\subseteq V.
\]
The same reweighting satisfies the stronger continuous inequality
\[
  (1-\eps)\sum_{e\in\mathcal E}N_e(x)
  \leq\sum_{e\in\mathcal E'}c'_eR_e(x)
  \leq(1+\eps)\sum_{e\in\mathcal E}N_e(x)
  \qquad(x\in\R^V),
\]
with zero-weight edges omitted.
\end{corollary}

\begin{proof}
Discard zero-weight hyperedges and apply Theorem~\ref{thm:main} to the
seminorms in~\eqref{eq:range-seminorm}.  Set
$\mathcal E'=\supp(w)$ and $c'_e=w_ec_e$.  The common kernel of the
$N_e$ consists exactly of vectors that are constant on each connected
component of the positive-weight $2$-section.  Its dimension is
$\kappa(H)$, so the effective dimension is $n-\kappa(H)$.  Evaluating
at indicator vectors gives the cut inequalities.
\end{proof}

\paragraph{Acknowledgments.}
The proof is built on the coefficient-space and partial-coloring
framework of Reis and Rothvoss~\cite{RR26}.  We thank the authors for
making their manuscript publicly available.

\bibliographystyle{alpha}
\bibliography{references}

\end{document}
